\origpage{166}{ROBERTSON AND SEYMOUR}
\textbf{(P4)} if $n_1, n_2$ are nodes bordering different cuffs, then $\beta(n_1) \cap \beta(n_2) = \emptyset$, and if $n_1$ is a node bordering a cuff $\Theta$ and $n_2 \in dom(\gamma)$ does not border $\Theta$, then $\gamma(n_2) \notin \beta(n_1)$;

\textbf{(P5)} if $c$ is an internal cell then $V(\alpha(c) \cap \alpha(-c)) \subseteq \gamma(\tilde{c})$;

\textbf{(P6)} if $c$ is a border cell with ends $n_1, n_2$, then
\begin{equation*}
V(\alpha(c) \cap \alpha(-c)) \subseteq \beta(n_1 +) \cup \beta(n_2 +);
\end{equation*}
\textbf{(P7)} if $n_1, n_2, n_3, n_4$ are nodes bordering the same cuff of $\Sigma$ and in order, then $\beta(n_1 +) \cap \beta(n_3 +) \subseteq \beta(n_2 +) \cup \beta(n_4 +)$.

(``In order,'' here and later in the paper, refers to the order of occurrence around the cuff. We see that the conclusion of \textbf{(P7)} holds trivially unless $n_1, n_2, n_3, n_4$ are all distinct.) For readers having trouble digesting this definition, it may help at this point to jump forward to (13.1) in order to better grasp what portrayals are.

Some remarks are in order. (1) Permitting $dom(\gamma) \ne N(\Gamma)$ is a helpful but artificial device to assist in removing vertices. We shall eventually be able to restrict ourselves to portrayals with $dom(\gamma) = N(\Gamma)$.

\hspace*{12pt}(2) Unfortunately, we need to permit $\gamma(n_1) \in \beta(n_2)$ for distinct nodes $n_1, n_2$ bordering the same cuff. Eventually this will be eliminated, but for the moment the reader is warned that it can happen.

\hspace*{12pt}(3) \textbf{(P6)} and \textbf{(P7)} tell us, in spirit, that for each cuff $\Theta$ the $\alpha(c)$'s for $c$ bordering $\Theta$ are arranged in a circle, each overlapping the next in $\beta(n +)$ (where $n$ is the common end). However, there is a possible degeneracy to beware of. If $c_1, c_2$ border $\Theta$, with ends $n_1, n_2$ and $n_2, n_3$, it is possible that $\beta(n_1 +) \cap \beta(n_3 +) \not\subseteq \beta(n_2 +)$; but any vertex in $\beta(n_1 +) \cap \beta(n_3 +) - \beta(n_2 +)$ lies in $\beta(n_4 +)$ for every $n_4 \ne n_2$ bordering $\Theta$. Eventually this degeneracy too will be eliminated.

The reason for interest in portrayals is that they provide a means of ``encoding'' the hypergraph $G$ by the painting $\Gamma$, labelling the cells of $\Gamma$ appropriately; and the product of the main theorem of [2] can be converted into a portrayal. The reason we are interested in portrayals of hypergraphs and not just of graphs is that, by applying [2] to the ``1-skeleton'' of a hypergraph, we obtain a portrayal of the hypergraph, as we shall see, and this will be important in proving Nash-Williams' ``immersions'' conjecture.

First, let us observe the following.

\stmt
(2.1) \textit{Given a portrayal $\pi = (\Sigma, \Gamma, \alpha, \beta, \gamma)$ of $G$, let $(I_1, I_2)$ be a partition of $C(\Gamma)$. For $i = 1, 2$, let $A_i = \bigcup (\alpha(c) : c \in I_i)$. Then $E(A_1 \cap A_2) = \emptyset$, and $V(A_1 \cap A_2) = \bigcup (\gamma(\tilde{c}_1) \cap \gamma(\tilde{c}_2) : c_1 \in I_1, c_2 \in I_2) \cup \bigcup (\beta(n) : \text{there exist border cells } c_1 \in I_1 \text{ and } c_2 \in I_2 \text{ with } n \in \tilde{c}_1 \cap \tilde{c}_2)$.}
